\documentclass[12pt]{article}
\usepackage{graphicx}
\usepackage{mathtools,amsfonts,amsthm}
\usepackage{lipsum}
\usepackage{enumitem}% http://ctan.org/pkg/enumitem

\begin{document}



\newpage
\section*{ABSTRACT}
Support for cancer patients from their caregivers is crucial in terms of logistical, emotional, and physical care. However, the difficulties they encounter in rural areas can raise their vulnerability to depression. This abstract seeks to present a succinct overview of the incidence, contributing factors, and effects of depression in rural carers of cancer patients.\\
A comprehensive literature review was conducted, focusing on studies published between 2009 and 2023. Relevant databases including PubMed, PsycINFO, and Google Scholar were searched using keywords related to depression, caregivers, cancer patients, and rural settings. The selection criteria included studies that specifically examined the prevalence, risk factors, and impact of depression on caregivers in rural areas.\\
According to the analysis, rural caregivers for cancer patients frequently worry about depression. Increased levels of depression are a result of the particular difficulties that caregivers in rural locations confront, including restricted access to healthcare resources, and financial strain. This risk is further increased by the stress of providing care, which includes the physical difficulties, emotional anguish, and strained social interactions.\\
Due to the difficulties they face, caregivers for cancer patients in rural locations are more likely to experience depression. For the creation of focused therapies and support programs, it is imperative to comprehend the particular elements that lead to depression in this population. Future studies should investigate practical ways to meet the mental health requirements of rural caregivers, such as expanding social support systems, increasing access to mental healthcare, and offering caregiver education and training.


\end{document}